\documentclass[12pt]{article}
\usepackage[utf8]{inputenc}
\usepackage[T1]{fontenc}
\usepackage[spanish,es-tabla]{babel}
\usepackage[left=1.54cm,right=1.54cm,top=1.5cm,bottom=1.8cm]{geometry}
\usepackage{amsmath,amssymb}
\usepackage{enumitem}
\usepackage{booktabs}
\usepackage{tikz}

\pagestyle{plain}
\setlength{\parindent}{0pt}
\setlength{\parskip}{0.4em}

\begin{document}
\shorthandoff{<>}
\large

% ---------- Encabezado ----------
\noindent
\begin{minipage}[t]{0.45\textwidth}
\raggedright
Mec\'anica de Fluidos CI3162\\
Primavera 2026
\end{minipage}%
\hfill
\begin{minipage}[t]{0.45\textwidth}
\raggedleft
Prof.\ Alberto de la Fuente Stranger\\
Auxs.\hspace{0.6cm}Josefina Moreno Faune\\
Sebasti\'an Fredes Ojeda
\end{minipage}

\vspace{3pt}
\hrule
\vspace{-5pt}

\begin{center}
{\Huge\bfseries Formulario --- Clase Auxiliar C1}\\[6pt]
Viernes 4 de Septiembre 2026
\end{center}

\vspace{14pt}

\noindent\textit{Contiene \'unicamente las herramientas de c\'atedra necesarias para
resolver los 4 problemas de la Clase Auxiliar C1. Donde una f\'ormula es producto de
una demostraci\'on vista en clases, se incluye su deducci\'on.}

\vspace{6pt}
\hrule
\vspace{10pt}

%======================================================================
\noindent{\LARGE\bfseries 1. Hidrost\'atica b\'asica}

Para un fluido en reposo bajo gravedad, $\nabla P=\rho\vec g$ ($z$ hacia arriba) da
\[
\frac{dP}{dz}=-\rho g \quad\Longrightarrow\quad \boxed{P(z)=P_0+\rho g\,h}
\]
con $h$ la profundidad desde el punto de presi\'on $P_0$. Es la base de los Problemas
2 y 3 (hidrost\'atica multicapa: la presi\'on se propaga sumando $\rho_i g\,\Delta h_i$
capa a capa).

%======================================================================
\noindent{\LARGE\bfseries 2. Ley de viscosidad de Newton}

Para un fluido newtoniano en corte simple, $\tau=\mu\dfrac{du}{dy}$, con $\tau$ el
esfuerzo de corte, $\mu$ la viscosidad din\'amica y $du/dy$ el gradiente de velocidad.
Es la ley constitutiva que se integra en la Herramienta 3.

%======================================================================
\noindent{\LARGE\bfseries 3. Viscos\'imetro de cilindros conc\'entricos (Couette anular)}

Dos cilindros conc\'entricos de radios $R_i<R_e$ y altura $H$, con lubricante de
viscosidad $\mu$ en el espacio anular. El exterior gira con velocidad angular $\omega$;
el interior permanece fijo. Para flujo de Stokes, axisim\'etrico y estacionario, la
velocidad tangencial satisface
\[
\frac{d}{dr}\left[\frac1r\frac{d(ru_\theta)}{dr}\right]=0
\quad\Longrightarrow\quad
u_\theta(r)=Ar+\frac{B}{r}
\]
Con las condiciones de borde $u_\theta(R_i)=0$ (cilindro interior fijo) y
$u_\theta(R_e)=\omega R_e$ (cilindro exterior), se obtiene
\[
A=\frac{\omega R_e^2}{R_e^2-R_i^2}\ ,\qquad B=-\frac{\omega R_e^2R_i^2}{R_e^2-R_i^2}
\]
El esfuerzo de corte es $\tau_{r\theta}=\mu r\dfrac{d}{dr}\!\left(\dfrac{u_\theta}{r}\right)=-\dfrac{2\mu B}{r^2}$,
y el torque por unidad de \'area en un cilindro de radio $r$ es
$T=2\pi H r^2\tau_{r\theta}=-4\pi H\mu B$, es decir
\[
\boxed{T=4\pi\mu H\omega\,\frac{R_e^2R_i^2}{R_e^2-R_i^2}}
\]
Usada en el Problema 1 (con $\omega=U/R_e$, ya que la cuerda se enrolla en el cilindro
exterior).

%======================================================================
\noindent{\LARGE\bfseries 4. Equilibrio de momentos de un cuerpo r\'igido}

Segunda ley de Newton para rotaci\'on en equilibrio (o traslaci\'on a velocidad
constante, sin aceleraci\'on), $\sum \vec\tau=0$ (o $\sum\vec F=0$). En el Problema 1,
el torque viscoso de la Herramienta 3 equilibra el torque que ejerce el peso del carro
a trav\'es de la cuerda (brazo $R_e$, componente del peso a lo largo del plano
$Mg\sin\alpha$):
\[
\boxed{T_{visc}=Mg\sin\alpha\cdot R_e}
\]

%======================================================================
\noindent{\LARGE\bfseries 5. Fuerza y torque de presi\'on sobre una superficie plana
sumergida}

Superficie plana de \'area $A$, con $y$ medida desde la superficie libre a lo largo del
plano (profundidad real $y\sin\theta$) y presi\'on de referencia $P_0$ en $y=0$:
$P(y)=P_0+\rho g\,y\sin\theta$.

\textbf{Fuerza resultante} ($y_G=\frac1A\iint y\,dA$, centroide):
\[
F=\iint P\,dA=P_0A+\rho g\sin\theta\,y_G A
\quad\Longrightarrow\quad
\boxed{F=\left(P_0+\rho g\,y_G\sin\theta\right)A}
\]

\textbf{Centro de presi\'on} $y_R$ (momento respecto de $y=0$, teorema de ejes
paralelos $\iint y^2dA=I_{xG}+Ay_G^2$; caso $P_0=0$):
\[
\boxed{y_R=y_G+\frac{I_{xG}}{y_G\,A}}
\]
Para un c\'irculo de di\'ametro $D$: $A=\pi D^2/4$, $I_{xG}=\pi D^4/64$. Usada en el
Problema 2 (compuerta circular con dos lados) para obtener $F_a,F_b$ (fuerzas) y
$b_a,b_b$ (brazos respecto de la r\'otula, $b=D/2+D^2/(16\,y_G)$), y el equilibrio de
momentos que da $F_{max}$:
\[
\boxed{F_{max}=\frac{F_b\,b_b-F_a\,b_a}{D}}
\]

%======================================================================
\noindent{\LARGE\bfseries 6. Fuerza vertical sobre una superficie curva (semiesfera)}

La componente vertical de la fuerza de presi\'on sobre una superficie curva es
$dF_z=P\,dA_h$ ($dA_h$ = proyecci\'on horizontal del elemento de \'area). Para una
compuerta semiesf\'erica de radio $R$, con presi\'on $P_5$ en su punto m\'as alto
(el borde, a profundidad cero) y l\'iquido de densidad $\rho_1$ llenando la calota,
integrando en coordenadas esf\'ericas ($z$ = profundidad bajo el borde, $r$ = radio
proyectado, $z=\sqrt{R^2-r^2}$):
\[
F_z=\int_0^R \left(P_5+\rho_1 g\sqrt{R^2-r^2}\right)2\pi r\,dr
= P_5\,\pi R^2+\rho_1 g\,\frac{2}{3}\pi R^3
\]
que es, como es de esperar, el resultado general (Herramienta anterior de fuerza sobre
superficie curva): presi\'on en el borde por el \'area proyectada, m\'as el peso del
volumen de la calota (semiesfera, $V=\frac23\pi R^3$). Descontando la presi\'on
atmosf\'erica que act\'ua por el lado exterior de la compuerta,
\[
\boxed{F=(P_5-P_{atm})\,\pi R^2-\rho_1 g\,\frac{2}{3}\pi R^3}
\]
es la fuerza neta que debe resistir la r\'otula (Problema 3).

%======================================================================
\noindent{\LARGE\bfseries 7. Ley de gases ideales}

\[
\boxed{PV=nRT}
\qquad\text{o, por unidad de masa,}\qquad
\boxed{P=\rho\frac{R}{m}T}
\]
con $R$ la constante universal de los gases ideales y $m$ la masa molar del gas.
Relaciona presi\'on, densidad (o volumen) y temperatura de cada gas en el Problema 4.

%======================================================================
\noindent{\LARGE\bfseries 8. Procesos termodin\'amicos}

Para una masa fija de gas ideal ($PV=nRT$ constante en $n$):
\begin{itemize}
\item \textbf{Isot\'ermico} ($T$ constante): $\boxed{P_iV_i=P_fV_f}$.
\item \textbf{Isoc\'orico} ($V$ constante): $\boxed{\dfrac{P_i}{T_i}=\dfrac{P_f}{T_f}}$.
\end{itemize}
Usados en la parte 2 del Problema 4 para relacionar los estados inicial y final de
cada gas.

%======================================================================
\noindent{\LARGE\bfseries 9. Presi\'on de Laplace (tensi\'on superficial) en una
burbuja de jab\'on}

Una interfaz curva esf\'erica de radio $r$ y tensi\'on superficial $\sigma$ sostiene
una diferencia de presi\'on $\Delta P=2\sigma/r$ entre su interior y exterior. Esto se
obtiene del equilibrio de fuerzas sobre un hemisferio de la interfaz: la fuerza neta de
presi\'on, $\Delta P\cdot\pi r^2$, es resistida por la tensi\'on superficial actuando a
lo largo del per\'imetro, $\sigma\cdot 2\pi r$:
\[
\Delta P\,\pi r^2=\sigma\, 2\pi r
\quad\Longrightarrow\quad
\Delta P=\frac{2\sigma}{r}
\]
Una \emph{burbuja de jab\'on} tiene dos interfaces (pel\'icula con dos superficies,
interior y exterior, ambas de radio $\approx d/2$ si el espesor $\epsilon\ll d$), por lo
que las dos ca\'idas de presi\'on se suman:
\[
\boxed{P_2-P_1=\frac{4\sigma}{d/2}=\frac{8\sigma}{d}}
\]
($P_2$ = presi\'on del gas interior a la burbuja, $P_1$ = presi\'on del gas exterior,
$d$ = di\'ametro interior de la burbuja). Junto con la Herramienta 7 (gas ideal, para
cada uno de los dos gases) permite despejar $d$ en el Problema 4.

\end{document}
